В условиях стремительного роста объемов информации классические подходы к управлению персональными знаниями становятся неэффективными. На практике все чаще применяются продвинутые решения, предоставляющие широкий функционал, в том числе генерацию, дополненную поиском (RAG). Однако широко используемый гибридный поиск имеет несколько недостатков. Более современный подход, использующий графы знаний, исправляет эти недочеты, но требует значительных трудозатрат на создание и поддержание, из-за чего с увеличением размера базы знаний пользователи сталкиваются с «узким местом», делающим ручное поддержание таких графов невозможным. Цель работы~--- разработка автоматизированного алгоритма для построения графов знаний с помощью семантического связывания заметок в экосистеме Obsidian, а также сравнение разработанного решения с существующими аналогами в сценариях использования RAG. Методология исследования основывается на методах обработки естественного языка, машинного обучения и информационного поиска. Эмпирической базой является разработанный валидационный набор данных на основе статей Википедии. Основной результат работы заключается в создании и оценке конвейера, автоматически находящего и типизирующего семантические связи. Разработанный алгоритм использует гибридный поиск, фильтрует результаты через Cross-Encoder, после чего передает их в LLM-типизатор, работающий на основе заданной онтологии. Показано, что предложенное решение устойчиво к шуму и демонстрирует качество связывания, сопоставимое с ручной разметкой. Также был создан и оценен усовершенствованный алгоритм GraphRAG для Obsidian, поддерживающий многошаговый обход графа и учет типов связей. Такой подход обеспечивает высокую полноту и точность извлечения информации, сопоставимые с продвинутыми корпоративными решениями.   
  
With the fast growth of information volumes, traditional approaches to personal knowledge management are becoming increasingly ineffective. In practice, advanced solutions providing a wide range of functionalities, including Retrieval-Augmented Generation (RAG), are gaining popularity. However, the widely used hybrid search method has several drawbacks. A more modern approach that utilizes knowledge graphs addresses these disadvantages but requires significant effort to create and maintain. As a consequence, as the knowledge base grows, users encounter a bottleneck that makes manual maintenance of such graphs impossible. The main goal of this work is to develop an automated algorithm for knowledge graph construction via the semantic linking of notes within the Obsidian ecosystem, as well as to compare the developed solution with existing alternatives in RAG use cases. The research methodology is based on natural language processing, machine learning, and information retrieval methods. The empirical foundation relies on a custom validation dataset constructed from Wikipedia articles. The main result of the study is the development and evaluation of a pipeline that automatically detects and classifies semantic relations. The algorithm uses a hybrid search strategy, filters the results via a Cross-Encoder, and then passes them to an LLM-based relation classifier that operates on a predefined ontology. The proposed solution is proven to be robust to noise and provides linking quality comparable to manual annotation. Furthermore, an improved GraphRAG algorithm for Obsidian was developed and evaluated. It supports multi-step graph traversal and incorporates relationship types. This approach achieves high recall and precision in information retrieval, comparable to advanced enterprise-grade solutions.  
